\documentclass[12pt,a4paper]{report}
% ─── Encoding & Language ───────────────────────────────────────────────────────
\usepackage[utf8]{inputenc}
\usepackage[T1]{fontenc}
\usepackage[romanian]{babel}

% ─── Page Layout ───────────────────────────────────────────────────────────────
\usepackage[a4paper, top=2.5cm, bottom=2.5cm, left=3cm, right=2cm]{geometry}
\usepackage{setspace}
\onehalfspacing

% ─── Fonts ─────────────────────────────────────────────────────────────────────
\usepackage{mathptmx}        % Times New Roman for text and math

% ─── Paragraph indent ──────────────────────────────────────────────────────────
\setlength{\parindent}{1.25cm}
\setlength{\parskip}{2pt}

% ─── Graphics & Colors ─────────────────────────────────────────────────────────
\usepackage{graphicx}
\usepackage{xcolor}
\usepackage{float}

% ─── Tables ────────────────────────────────────────────────────────────────────
\usepackage{booktabs}
\usepackage{longtable}
\usepackage{array}
\usepackage{tabularx}

% ─── Code Listings ─────────────────────────────────────────────────────────────
\usepackage{listings}
\usepackage{lstautogobble}

\definecolor{codebg}{RGB}{245,245,245}
\definecolor{codeframe}{RGB}{200,200,200}
\definecolor{keyword}{RGB}{0,0,180}
\definecolor{string}{RGB}{163,21,21}
\definecolor{comment}{RGB}{0,128,0}
\definecolor{placeholderbg}{RGB}{230,228,224}

\lstset{
  backgroundcolor=\color{codebg},
  frame=single,
  rulecolor=\color{codeframe},
  basicstyle=\ttfamily\small,
  keywordstyle=\color{keyword}\bfseries,
  stringstyle=\color{string},
  commentstyle=\color{comment}\itshape,
  breaklines=true,
  breakatwhitespace=true,
  showstringspaces=false,
  numbers=left,
  numberstyle=\tiny\color{gray},
  numbersep=6pt,
  tabsize=4,
  autogobble=true,
  captionpos=b,
  aboveskip=6pt,
  belowskip=6pt,
}

\lstdefinelanguage{TypeScript}{
  keywords={const, let, var, function, return, import, export, from, type,
            interface, class, extends, implements, new, this, if, else,
            for, while, switch, case, break, default, async, await, string,
            number, boolean, void, null, undefined, true, false, error},
  morestring=[b]',
  morestring=[b]",
  morestring=[b]`,
  morecomment=[l]{//},
  morecomment=[s]{/*}{*/},
  sensitive=true
}

\lstdefinelanguage{Svelte}{
  keywords={script, style, div, h1, h2, h3, p, a, button, nav, img, section,
            article, header, footer, span, ul, li, svelte, head,
            const, let, var, function, return, import, export, from, if,
            else, each, onMount, bind, this},
  morestring=[b]',
  morestring=[b]",
  morecomment=[l]{//},
  morecomment=[s]{/*}{*/},
  sensitive=true
}

% ─── Image placeholder ─────────────────────────────────────────────────────────
\newcommand{\imgplaceholder}[3][6cm]{%
  \begin{figure}[H]
    \centering
    \fcolorbox{codeframe}{placeholderbg}{%
      \parbox[c][#1][c]{0.88\textwidth}{%
        \centering\color{gray}\itshape #2%
      }%
    }
    \caption{#2}
    \label{#3}
  \end{figure}
}

% ─── Hyperlinks ────────────────────────────────────────────────────────────────
\usepackage{hyperref}
\hypersetup{
  colorlinks=true,
  linkcolor=black,
  urlcolor=blue,
  citecolor=black,
  pdftitle={Lucrare de Atestat -- Frank Lloyd Wright -- Arhitectura Organica},
  pdfauthor={Daniliuc Mirela Iasmina},
}

% ─── Headers & Footers ─────────────────────────────────────────────────────────
\usepackage{fancyhdr}
\pagestyle{fancy}
\fancyhf{}
\fancyfoot[C]{\thepage}
\renewcommand{\headrulewidth}{0pt}

% ─── Chapter / Section Formatting ──────────────────────────────────────────────
\usepackage{titlesec}
\titleformat{\chapter}[block]
  {\normalfont\Large\bfseries}{\thechapter.}{0.8em}{}
\titlespacing*{\chapter}{0pt}{14pt}{8pt}

\titleformat{\section}[block]
  {\normalfont\large\bfseries}{\thesection}{0.8em}{}
\titlespacing*{\section}{0pt}{10pt}{5pt}

\titleformat{\subsection}[block]
  {\normalfont\normalsize\bfseries}{\thesubsection}{0.8em}{}
\titlespacing*{\subsection}{0pt}{8pt}{4pt}

% ─── List Spacing ──────────────────────────────────────────────────────────────
\usepackage{enumitem}
\setlist{
  itemsep=1pt,
  parsep=0pt,
  topsep=3pt,
  partopsep=0pt,
}
\setlist[itemize]{leftmargin=1.5cm}
\setlist[enumerate]{leftmargin=1.5cm}

% ─── Miscellaneous ─────────────────────────────────────────────────────────────
\usepackage{microtype}
\usepackage{caption}
\captionsetup{font=small, labelfont=bf, skip=4pt}

% ═══════════════════════════════════════════════════════════════════════════════
\begin{document}

% ─── COVER PAGE ────────────────────────────────────────────────────────────────
\begin{titlepage}
  \centering
  {\Large\bfseries Colegiul Național „A.T. Laurian" Botoșani\par}
  \vspace{0.2cm}
  {\large Specializarea: Matematică-informatică, intensiv engleză\par}
  \vspace{1.5cm}
  \vfill
  {\huge\bfseries Lucrare de atestat profesional\\[0.3cm] la informatică\par}
  \vfill

  {\large\bfseries Profesor coordonator:}\\[0.2cm]
  {\large Giocaș Carmen Afrodita\par}
  \vspace{1cm}
  {\large\bfseries Candidat:}\\[0.2cm]
  {\large Daniliuc Mirela Iasmina\par}
  \vspace{1.5cm}
  {\large Mai, 2026}
\end{titlepage}


% ─── PROJECT TITLE PAGE ────────────────────────────────────────────────────────
\thispagestyle{empty}
\vspace*{\fill}
\begin{center}
  {\large\bfseries Tema Proiectului:\par}
  \vspace{0.5cm}
  {\LARGE\bfseries Frank Lloyd Wright\\[0.3cm]
   Arhitectură Organică\par}
  \vspace{0.4cm}
  {\large\itshape Prezentare web a operei arhitecturale a lui Frank Lloyd Wright\par}
\end{center}
\vspace*{\fill}
\newpage

% ─── TABLE OF CONTENTS ─────────────────────────────────────────────────────────
\tableofcontents
\newpage

% ═══════════════════════════════════════════════════════════════════════════════
%  CHAPTER 1 — ARGUMENT
% ═══════════════════════════════════════════════════════════════════════════════
\chapter{Argument}

Am ales ca temă pentru lucrarea mea de atestat realizarea unui site web dedicat
operei arhitecturale a lui \textbf{Frank Lloyd Wright} (1867--1959), cel mai
influent arhitect american al secolului al XX-lea. Proiectul prezintă, în limba
română, opt dintre cele mai reprezentative lucrări ale sale, de la Fallingwater
și Muzeul Guggenheim până la casele Usonian accesibile clasei de mijloc.

Motivul alegerii acestei teme îl constituie intersecția dintre două domenii care
mă pasionează: \textbf{informatica} și \textbf{designul vizual}. Frank Lloyd
Wright a revoluționat arhitectura prin conceptul de \emph{arhitectură organică}
--- ideea că o clădire trebuie să crească armonios din sitorul, materialele și
nevoile utilizatorilor ei. Am găsit în această filozofie o paralelă puternică
cu principiile bunului design de interfețe: un site bun, la fel ca o clădire
bună, trebuie să fie funcțional, coerent și să comunice conținutul fără efort.

Realizarea proiectului a constituit o oportunitate de a aplica în practică
concepte din domeniul dezvoltării web moderne: rutare bazată pe sistem de fișiere,
generare statică a paginilor, animații CSS controlate prin JavaScript nativ și
exportul unui site complet funcțional ca fișiere HTML statice, deschizabile
direct din browser fără a necesita un server web.

% ═══════════════════════════════════════════════════════════════════════════════
%  CHAPTER 2 — INTRODUCERE
% ═══════════════════════════════════════════════════════════════════════════════
\chapter{Introducere}

Lucrarea de față prezintă site-ul web \textbf{Frank Lloyd Wright --- Arhitectură
Organică}, o prezentare digitală a operei arhitecturale a lui Frank Lloyd Wright,
realizată în limba română.

\section{Paginile aplicației}

Din meniul principal al aplicației, utilizatorul poate naviga către următoarele
secțiuni:

\begin{itemize}
  \item \textbf{Acasă} --- Pagina principală cu secțiunea hero, introducere
        biografică, proiecte selectate, citat filozofic și principii de design
  \item \textbf{Proiecte} --- Lista completă a celor opt proiecte arhitecturale,
        cu imagini interactive și detalii tehnice
  \item \textbf{Proiect individual} --- Pagina detaliată a fiecărei lucrări, cu
        galerie foto, descriere aprofundată și navigare între proiecte
  \item \textbf{Despre} --- Biografie, linie de timp și filosofia arhitecturală
        a lui Frank Lloyd Wright
\end{itemize}

\section{Tehnologiile utilizate}

Aplicația a fost realizată utilizând:

\begin{itemize}
  \item \textbf{SvelteKit 2.x} cu \textbf{Svelte 5} --- framework pentru
        construirea interfețelor web, cu sistemul reactiv de tip Runes
  \item \textbf{TypeScript} --- pentru tipizarea statică a datelor și a
        logicii componentelor
  \item \textbf{CSS scoped} --- stilizare izolată per componentă, fără
        framework extern, cu variabile CSS globale pentru sistemul de design
  \item \textbf{@sveltejs/adapter-static} --- pentru generarea unui site
        complet static, exportabil ca fișiere HTML/CSS
  \item \textbf{Vite 8.x} --- instrument de build și server de dezvoltare
\end{itemize}

\imgplaceholder[5.5cm]{Pagina principală a site-ului Frank Lloyd Wright}{fig:home}

% ═══════════════════════════════════════════════════════════════════════════════
%  CHAPTER 3 — GENERALITĂȚI DESPRE TEHNOLOGIILE FOLOSITE
% ═══════════════════════════════════════════════════════════════════════════════
\chapter{Generalități despre tehnologiile folosite}

\section{Limbajul TypeScript}

\textbf{TypeScript} este un limbaj de programare dezvoltat de \textbf{Microsoft},
lansat în 2012, care extinde JavaScript prin adăugarea \textbf{tipizării statice}.
Codul TypeScript este transpilat în JavaScript standard înainte de rulare, iar
erorile de tip sunt detectate la momentul compilării, nu în producție.

Principalele avantaje ale TypeScript:

\begin{itemize}
  \item \textbf{Detectarea erorilor la compilare:} Greșelile de tip sunt
        identificate înainte de rularea codului.
  \item \textbf{Autocompletare inteligentă:} Editoarele de cod pot oferi
        sugestii precise bazate pe tipuri.
  \item \textbf{Interfețe și tipuri:} Permit definirea clară a structurii
        datelor complexe.
  \item \textbf{Compatibilitate cu JavaScript:} Orice cod JavaScript valid
        este și cod TypeScript valid.
\end{itemize}

În cadrul proiectului, TypeScript este utilizat pentru definirea modelului de
date al proiectelor arhitecturale (fișierul \texttt{src/lib/data/projects.ts})
și pentru logica de navigare și animații din componente.

\begin{lstlisting}[language=TypeScript,
  caption={Interfața \texttt{Project} din \texttt{src/lib/data/projects.ts}}]
export interface Project {
    slug: string;
    title: string;
    year: number;
    location: string;
    style: string;
    client: string;
    material: string;
    intro: string;
    description: string;
    images: {
        hero: string;
        exterior: string;
        exterior2?: string;
        interior: string;
        interior2?: string;
        withWright?: string;
    };
}
\end{lstlisting}

\section{Framework-ul SvelteKit}

\textbf{Svelte} este un framework UI creat de \textbf{Rich Harris}, lansat în
2016. Spre deosebire de React sau Vue, Svelte \textbf{compilează componentele
la momentul construirii} (\textit{build time}), generând JavaScript optimizat
fără a necesita un DOM virtual la runtime. Rezultatul este o aplicație mai
rapidă și cu un pachet mai mic.

\textbf{SvelteKit} (versiunea 2.x) este framework-ul full-stack construit pe
baza Svelte. Oferă:

\begin{itemize}
  \item \textbf{Rutare bazată pe sistemul de fișiere:} Fiecare fișier
        \texttt{+page.svelte} dintr-un director devine o pagină la acea rută.
  \item \textbf{Generare statică (SSG):} Prin \texttt{adapter-static}, toate
        paginile sunt pre-randate la build time ca fișiere HTML pure.
  \item \textbf{Layout-uri partajate:} Fișierul \texttt{+layout.svelte}
        definește structura comună (bară de navigare, subsol).
  \item \textbf{Rute dinamice:} Directoarele \texttt{[slug]} generează câte o
        pagină pentru fiecare parametru definit în funcția \texttt{entries()}.
  \item \textbf{Prerendare completă:} Exportul \texttt{prerender = true}
        instruiește SvelteKit să genereze HTML pentru toate paginile la build.
\end{itemize}

\textbf{Svelte 5} (utilizat în proiect) introduce conceptul de \textbf{Runes}
--- un nou sistem reactiv bazat pe \texttt{\$state}, \texttt{\$derived} și
\texttt{\$props} pentru gestionarea stării componentelor, înlocuind sintaxa
anterioară bazată pe \texttt{let} reactiv și \texttt{store}.

\begin{lstlisting}[language=Svelte,
  caption={Exemplu de Runes Svelte 5 în componenta Navbar}]
<script lang="ts">
    import { page } from '$app/stores';

    const links = [
        { href: '/', label: 'acasa' },
        { href: '/proiecte', label: 'proiecte' },
        { href: '/despre', label: 'despre' }
    ];

    let scrolled = $state(false);

    function onScroll() {
        scrolled = window.scrollY > 40;
    }
</script>

<svelte:window on:scroll={onScroll} />

<nav class:scrolled>
    <!-- ... -->
</nav>
\end{lstlisting}

\section{CSS Scoped și Variabile CSS}

Proiectul nu utilizează un framework CSS extern (precum TailwindCSS sau
Bootstrap), ci \textbf{CSS scoped}, stiluri izolate per componentă generate
automat de compilatorul Svelte. Fiecare componentă primește un sufix unic de
clasă (de exemplu \texttt{.svelte-rfuq4y}), prevenind conflictele de stiluri
între componente diferite.

Sistemul de design global este definit prin \textbf{variabile CSS} în
\texttt{:root}, în fișierul \texttt{src/app.css}:

\begin{lstlisting}[language=bash,
  caption={Variabilele CSS globale ale sistemului de design}]
:root {
  --bg: #F5F2ED;          /* fundal cald, aproape alb */
  --bg-alt: #EDEAE3;
  --ink: #0C0B09;         /* text principal */
  --ink-secondary: #3E3A35;
  --muted: #7A746D;
  --gold: #A87840;        /* culoarea de accent */
  --gold-light: #C4975A;
  --border: #DAD4CB;
  --dark-bg: #181512;     /* fundal subsol */
  --dark-text: #F5F2ED;

  --font-display: 'Space Grotesk', sans-serif;
  --font-body: 'DM Sans', sans-serif;

  --nav-height: 68px;
  --pad-x: clamp(24px, 5.5vw, 88px);
  --pad-section: clamp(80px, 11vw, 148px);
}
\end{lstlisting}

Paleta de culori evocă materialele naturale asociate lui Wright: crem-lut
pentru fundal, ocru-auriu pentru accent și aproape-negru organic pentru text.

\section{Generarea Statică --- adapter-static}

\textbf{@sveltejs/adapter-static} este adaptorul oficial SvelteKit pentru
generarea site-urilor complet statice. Cu această configurație, procesul de
build parcurge fiecare rută definită în aplicație, execută logica de
server-side rendering pentru acea rută și salvează rezultatul HTML în fișiere
pe disc, în directorul \texttt{build/}.

Configurația proiectului (\texttt{svelte.config.js}):

\begin{lstlisting}[language=bash,
  caption={Configurația SvelteKit pentru generare statică}]
import adapter from '@sveltejs/adapter-static';

export default {
  kit: {
    adapter: adapter({
      pages: 'build',
      assets: 'build',
      fallback: undefined
    }),
    paths: {
      relative: true    // Cai relative in loc de absolute
    },
    prerender: {
      entries: ['*']    // Genereaza toate paginile
    }
  }
};
\end{lstlisting}

Opțiunea \texttt{paths.relative: true} face ca toate referințele la resurse
(CSS, imagini, linkuri interne) să fie relative față de locul fișierului HTML,
ceea ce permite deschiderea site-ului direct din sistemul de fișiere, fără
server web.

Directiva \texttt{export const csr = false} din \texttt{+layout.ts} dezactivează
complet runtime-ul JavaScript al SvelteKit în browser: paginile sunt HTML pur,
fără hidratare. Animațiile și interactivitățile sunt gestionate prin JavaScript
nativ (vanilla JS), inclus direct în \texttt{app.html}.

\section{Vite --- Instrumentul de build}

\textbf{Vite} (versiunea 8.x) este instrumentul de build și serverul de
dezvoltare al proiectului. Folosește \textbf{ES Modules native} în browser
pentru pornire instantanee în modul de dezvoltare și \textbf{Rollup} pentru
bundling la producție.

Plugin-uri Vite utilizate:

\begin{itemize}
  \item \texttt{@sveltejs/vite-plugin-svelte} --- compilarea componentelor
        Svelte și integrarea cu SvelteKit
\end{itemize}

% ═══════════════════════════════════════════════════════════════════════════════
%  CHAPTER 4 — CERINTE HARDWARE SI SOFTWARE
% ═══════════════════════════════════════════════════════════════════════════════
\chapter{Cerințe hardware și software}

\section{Cerințe hardware pentru dezvoltare}

\begin{itemize}
  \item Procesor la minimum 1 GHz (recomandat 2 GHz sau mai mult)
  \item Minimum 2 GB RAM (recomandat 8 GB)
  \item Minimum 300 MB spațiu liber pe disc
  \item Conexiune la internet (pentru instalarea dependențelor npm și
        încărcarea fonturilor Google Fonts în browser)
\end{itemize}

\section{Software necesar pentru dezvoltare}

\begin{itemize}
  \item \textbf{Node.js} versiunea 18 sau mai recentă
  \item \textbf{npm} (versiunea 9 sau mai recentă) --- managerul de pachete
  \item \textbf{Git} --- pentru versionarea codului sursă
  \item Un editor de cod, recomandat \textbf{Visual Studio Code} cu extensia
        oficială Svelte (\texttt{svelte.svelte-vscode})
\end{itemize}

\section{Vizualizarea site-ului static (utilizator final)}

Spre deosebire de site-urile web tradiționale, versiunea exportată a
proiectului poate fi deschisă direct din sistemul de fișiere, fără a instala
sau porni vreun server web:

\begin{itemize}
  \item Orice browser web modern: Google Chrome (recomandat), Mozilla Firefox,
        Microsoft Edge, Safari
  \item Nu este necesar niciun software instalat pe calculatorul utilizatorului
        în afara unui browser
  \item Se deschide fișierul \texttt{build/index.html} direct din manager-ul
        de fișiere sau din linia de comandă
\end{itemize}

\section{Comenzi de dezvoltare}

\begin{lstlisting}[language=bash, caption={Comenzile npm disponibile în proiect}]
npm install          # Instalarea dependentelor
npm run dev          # Pornirea serverului de dezvoltare (Vite)
npm run build        # Build complet: vite build + relativizare linkuri
npm run preview      # Previzualizarea build-ului pe un server local
npm run check        # Verificare tipuri TypeScript si Svelte
npm run format       # Formatare automata a codului sursa (Prettier)
\end{lstlisting}

Comanda \texttt{npm run build} execută două etape succesive: mai întâi
\texttt{vite build} generează fișierele HTML statice în directorul
\texttt{build/}, apoi scriptul \texttt{node scripts/relativize.js} rescrie
toate linkurile interne din formatul absolut (\texttt{href="/proiecte"}) în
formatul relativ (\texttt{href="./proiecte.html"}), necesar pentru protocolul
\texttt{file://}.

% ═══════════════════════════════════════════════════════════════════════════════
%  CHAPTER 5 — STRUCTURA SI CONTINUTUL PROIECTULUI
% ═══════════════════════════════════════════════════════════════════════════════
\chapter{Structura și conținutul proiectului}

\section{Harta site-ului}

\begin{table}[H]
  \centering
  \caption{Rutele aplicației și fișierele corespunzătoare}
  \label{tab:routes}
  \begin{tabularx}{\textwidth}{>{\ttfamily}l >{\ttfamily}l X}
    \toprule
    \textbf{Rută} & \textbf{Fișier pagină} & \textbf{Descriere} \\
    \midrule
    /                    & +page.svelte              & Pagina principală (hero, intro, proiecte selectate) \\
    /proiecte            & proiecte/+page.svelte     & Lista completă a celor 8 proiecte \\
    /proiecte/[slug]     & proiecte/[slug]/+page.svelte & Pagina detaliată a unui proiect \\
    /despre              & despre/+page.svelte       & Biografie, linie de timp, citate \\
    \bottomrule
  \end{tabularx}
\end{table}

La build, ruta dinamică \texttt{/proiecte/[slug]} generează 8 fișiere HTML
statice: \texttt{fallingwater.html}, \texttt{guggenheim.html},
\texttt{taliesin-west.html}, \texttt{robie-house.html},
\texttt{unity-temple.html}, \texttt{hollyhock-house.html},
\texttt{jacobs-house.html} și \texttt{taliesin.html}.

\section{Structura directoarelor}

\begin{lstlisting}[language=bash,
  caption={Structura proiectului franklloyd/}]
franklloyd/
|-- src/
|   |-- routes/
|   |   |-- +page.svelte           # Pagina principala
|   |   |-- +layout.svelte         # Layout global (Navbar + subsol)
|   |   |-- +layout.ts             # prerender=true, csr=false
|   |   |-- proiecte/
|   |   |   |-- +page.svelte       # Lista proiecte cu toggle imagine
|   |   |   |-- [slug]/
|   |   |   |   |-- +page.svelte   # Detaliu proiect (galerie, descriere)
|   |   |   |   |-- +page.ts       # entries() + load() cu date proiect
|   |   |-- despre/
|   |   |   |-- +page.svelte       # Biografie si linie de timp
|   |-- lib/
|   |   |-- components/
|   |   |   |-- Navbar.svelte      # Bara de navigare fixa
|   |   |   |-- ImageTilt.svelte   # Componenta imagine cu efect tilt
|   |   |-- data/
|   |   |   |-- projects.ts        # Date pentru cele 8 proiecte
|   |-- app.css                    # Stiluri globale si variabile CSS
|   |-- app.html                   # Template HTML de baza cu JS nativ
|-- scripts/
|   |-- relativize.js              # Post-build: rescrie linkurile ca relative
|-- static/
|   |-- placeholder.png            # Imagine placeholder pentru proiecte
|-- svelte.config.js               # Configuratie adapter-static
|-- vite.config.ts
|-- package.json
\end{lstlisting}

\section{Datele proiectelor arhitecturale}

Toate datele despre proiectele arhitecturale sunt stocate în fișierul
\texttt{src/lib/data/projects.ts}, ca un array de obiecte TypeScript. Fiecare
obiect respectă interfața \texttt{Project} și conține:

\begin{itemize}
  \item \textbf{slug} --- identificatorul unic folosit în URL (ex: \texttt{fallingwater})
  \item \textbf{title, year, location, style} --- date de identificare
  \item \textbf{client, material} --- informații tehnice și istorice
  \item \textbf{intro} --- rezumat scurt (afișat în liste și carduri)
  \item \textbf{description} --- descriere detaliată (afișată pe pagina individuală)
  \item \textbf{images} --- obiect cu URL-uri pentru imaginile hero, exterior și interior
\end{itemize}

Proiectele documentate sunt:

\begin{table}[H]
  \centering
  \caption{Proiectele arhitecturale prezentate în aplicație}
  \label{tab:projects}
  \begin{tabularx}{\textwidth}{l l >{\ttfamily}l X}
    \toprule
    \textbf{Nr.} & \textbf{An} & \textbf{Slug} & \textbf{Titlu} \\
    \midrule
    1 & 1910 & robie-house     & Casa Frederick C. Robie \\
    2 & 1908 & unity-temple    & Unity Temple \\
    3 & 1911 & taliesin        & Taliesin \\
    4 & 1921 & hollyhock-house & Hollyhock House \\
    5 & 1935 & fallingwater    & Fallingwater \\
    6 & 1937 & taliesin-west   & Taliesin West \\
    7 & 1937 & jacobs-house    & Casa Herbert Jacobs I \\
    8 & 1959 & guggenheim      & Muzeul Solomon R. Guggenheim \\
    \bottomrule
  \end{tabularx}
\end{table}

\section{Componenta Navbar}

Bara de navigare este implementată în \texttt{src/lib/components/Navbar.svelte}
și este inclusă global prin \texttt{+layout.svelte}. Are poziție fixă în
partea superioară a paginii și devine opacă (fundal solid) când utilizatorul
derulează pagina mai jos de 40 de pixeli, pentru a asigura lizibilitatea
linkurilor pe orice fond.

Legătura activă este detectată prin compararea \texttt{\$page.url.pathname}
cu URL-ul fiecărui element de navigare --- la prerendare, această valoare este
cunoscută, deci clasa \texttt{active} este baked in HTML la build time,
fără a necesita JavaScript în browser.

\begin{lstlisting}[language=Svelte,
  caption={Marcarea legăturii active în Navbar.svelte}]
{#each links as { href, label }}
    <li>
        <a {href}
           class:active={
               $page.url.pathname === href ||
               ($page.url.pathname.startsWith(href)
                && href !== '/')
           }>
            {label}
        </a>
    </li>
{/each}
\end{lstlisting}

\imgplaceholder[5cm]{Bara de navigare în stare normală și cu fundal activat la scroll}{fig:navbar}

\section{Componenta ImageTilt}

\texttt{src/lib/components/ImageTilt.svelte} implementează efectul de
perspectivă 3D la hover --- imaginea se încline ușor în direcția cursorului
mouse-ului, creând iluzia de adâncime. Efectul este realizat prin interpolarea
liniară (\textit{lerp}) a unghiurilor de rotație, animat printr-un loop
\texttt{requestAnimationFrame}.

\begin{lstlisting}[language=Svelte,
  caption={Logica de animație lerp din ImageTilt.svelte}]
function loop() {
    cx = lerp(cx, tx, 0.09);
    cy = lerp(cy, ty, 0.09);
    if (!hovering && Math.abs(cx) < 0.004
                  && Math.abs(cy) < 0.004) {
        cx = 0; cy = 0;
        imageEl.style.transform = '';
        running = false;
        return;
    }
    imageEl.style.transform =
        'perspective(900px) ' +
        'rotateY(' + cx + 'deg) ' +
        'rotateX(' + (-cy) + 'deg) ' +
        'scale(1.04)';
    raf = requestAnimationFrame(loop);
}
\end{lstlisting}

Componenta acceptă proprietățile: \texttt{src}, \texttt{alt},
\texttt{aspectRatio}, \texttt{intensity} și \texttt{overlay}. Intensitatea
efectului este stocată și în atributul HTML \texttt{data-intensity}, pentru
a putea fi citit de JavaScript nativ la redeschiderea paginii statice.

\imgplaceholder[5cm]{Efectul ImageTilt activ pe un card de proiect}{fig:imagetilt}

\section{Animațiile de tip scroll reveal}

Elementele de conținut sunt inițial invizibile (\texttt{opacity: 0},
\texttt{transform: translateY(36px)}) și devin vizibile la intrarea în
viewport prin clasa CSS \texttt{.visible}:

\begin{lstlisting}[language=bash,
  caption={Stilul de bază pentru animația reveal în app.css}]
.reveal {
  opacity: 0;
  transform: translateY(36px);
  transition:
    opacity 0.9s cubic-bezier(0.22, 1, 0.36, 1),
    transform 0.9s cubic-bezier(0.22, 1, 0.36, 1);
}
.reveal.visible {
  opacity: 1;
  transform: translateY(0);
}
.reveal-delay-1 { transition-delay: 0.1s; }
.reveal-delay-2 { transition-delay: 0.2s; }
\end{lstlisting}

Deoarece site-ul exportat dezactivează runtime-ul SvelteKit
(\texttt{csr = false}), logica \texttt{IntersectionObserver} este implementată
ca JavaScript nativ în \texttt{app.html}, executat la fiecare încărcare de
pagină:

\begin{lstlisting}[language=bash,
  caption={IntersectionObserver nativ în app.html (vanilla JS)}]
(function () {
    var observer = new IntersectionObserver(function (entries) {
        entries.forEach(function (e) {
            if (e.isIntersecting) e.target.classList.add('visible');
        });
    }, { threshold: 0.1 });
    document.querySelectorAll('.reveal').forEach(function (el) {
        observer.observe(el);
    });
})();
\end{lstlisting}

\section{Pagina principală}

Pagina principală este structurată în șase secțiuni:

\begin{enumerate}
  \item \textbf{Hero} --- Imagine full-screen cu titlul animat al arhitectului,
        cifre statistice (532 proiecte, 70+ ani de activitate) și un buton de
        scroll, animate la intrarea pe pagină printr-o clasă \texttt{visible}
        aplicată cu un \texttt{setTimeout} de 100 ms
  \item \textbf{Intro} --- Citat celebru și biografie scurtă, cu link spre
        pagina \emph{Despre}
  \item \textbf{Featured} --- Trei carduri cu proiectele cele mai reprezentative
        (Fallingwater, Guggenheim, Taliesin West), fiecare utilizând componenta
        \texttt{ImageTilt}
  \item \textbf{Philosophy band} --- Bandă cu citat filozofic suprapus pe o
        imagine de fundal
  \item \textbf{Principii fundamentale} --- Patru principii ale arhitecturii
        organice (orizontalitate, spațiu deschis, lumina naturală, materiale
        locale), prezentate în grilă
\end{enumerate}

\imgplaceholder[5.5cm]{Secțiunea hero a paginii principale, cu titlul animat}{fig:homepage-hero}

\section{Pagina Proiecte}

Pagina \texttt{/proiecte} prezintă toate cele opt proiecte în format listă, câte
unul pe rând. Fiecare intrare în listă conține:

\begin{itemize}
  \item Un număr de ordine și metadate (an, locație)
  \item Două imagini suprapuse (exterior/interior) cu posibilitate de comutare
        printr-un buton thumbnail; comutarea este gestionată de JavaScript nativ
  \item Efect de perspectivă tilt la hover, aplicat pe blocul de imagini
  \item Titlul, stilul arhitectural, o descriere introductivă și detalii despre
        client și materiale
  \item Un link \emph{descoperă proiectul} spre pagina individuală
\end{itemize}

\imgplaceholder[6cm]{Pagina Proiecte --- un rând din lista de lucrări arhitecturale}{fig:proiecte}

\section{Pagina individuală de proiect}

Fiecare proiect individual (\texttt{/proiecte/[slug]}) are o structură proprie:

\begin{enumerate}
  \item \textbf{Hero section} --- Titlul proiectului animat cuvânt cu cuvânt,
        stilul arhitectural și un tabel cu datele principale (an, locație,
        client, materiale)
  \item \textbf{Galerie exterior} --- Două imagini cu exteriorul clădirii,
        principale și secundare
  \item \textbf{Galerie interior} --- O imagine mare și una mai mică a
        interiorului
  \item \textbf{Descriere} --- Textul detaliat al proiectului (300--400 cuvinte)
  \item \textbf{Navigare} --- Săgeți spre proiectul precedent și următor
\end{enumerate}

Funcția \texttt{entries()} din \texttt{+page.ts} informează SvelteKit despre
toate valorile posibile ale parametrului \texttt{[slug]}, necesare pentru
prerendarea statică:

\begin{lstlisting}[language=TypeScript,
  caption={Funcția entries() pentru generarea statică a paginilor dinamice}]
import { projects } from '$lib/data/projects';
import { error } from '@sveltejs/kit';

export const entries = () =>
    projects.map((p) => ({ slug: p.slug }));

export const load = ({ params }) => {
    const project = projects.find(
        (p) => p.slug === params.slug
    );
    if (!project) error(404, 'Proiect negasit');

    const idx = projects.findIndex(
        (p) => p.slug === params.slug
    );
    return {
        project,
        prev: idx > 0 ? projects[idx - 1] : null,
        next: idx < projects.length - 1
            ? projects[idx + 1] : null
    };
};
\end{lstlisting}

\imgplaceholder[6cm]{Pagina individuală a proiectului Fallingwater}{fig:proiect-detail}

\section{Pagina Despre}

Pagina \texttt{/despre} prezintă biografia lui Frank Lloyd Wright prin trei
elemente principale:

\begin{itemize}
  \item \textbf{Linie de timp} --- Zece momente-cheie din viața arhitectului,
        de la nașterea din 1867 până la deschiderea Muzeului Guggenheim în 1959
  \item \textbf{Galerie portret} --- Imagini ale arhitectului și ale unor
        proiecte semnificative, cu efect \texttt{ImageTilt}
  \item \textbf{Citate filozofice} --- Trei citate celebre ale lui Wright,
        rotite automat la un interval de 5 secunde printr-un timer JavaScript
\end{itemize}

\imgplaceholder[5cm]{Pagina Despre --- secțiunea cu linia de timp}{fig:despre}

\section{Scriptul de relativizare a linkurilor}

O provocare tehnică a proiectului a fost compatibilizarea site-ului cu
protocolul \texttt{file://}: browserele moderne blochează navigarea spre URL-uri
absolute (\texttt{href="/proiecte"}) când un fișier HTML este deschis din
sistemul de fișiere local.

Scriptul \texttt{scripts/relativize.js} rezolvă această problemă prin
postprocesarea HTML-ului generat de SvelteKit:

\begin{lstlisting}[language=bash,
  caption={Scriptul relativize.js --- reescrierea linkurilor după build}]
const files = findHtmlFiles(buildDir);

for (const file of files) {
    const dir = dirname(file);
    const content = readFileSync(file, 'utf8')
        .replace(/href="(\/[^"#?]*)"/g, (_, href) => {
            // Converteste /proiecte -> ./proiecte.html
            const target = href === '/'
                ? 'index.html'
                : href.replace(/^\//, '') + '.html';
            let rel = relative(dir,
                join(buildDir, target));
            if (!rel.startsWith('.'))
                rel = './' + rel;
            return `href="${rel}"`;
        });
    writeFileSync(file, content);
}
\end{lstlisting}

Scriptul parcurge recursiv directorul \texttt{build/}, găsește toate fișierele
\texttt{.html} (excluzând directorul \texttt{\_app} cu resursele statice),
și rescrie fiecare atribut \texttt{href} cu cale absolută în calea relativă
corespunzătoare față de locația acelui fișier HTML.

\section{JavaScript nativ pentru interactivitate}

Deoarece runtime-ul SvelteKit este dezactivat în versiunea statică
(\texttt{csr = false}), toate interactivitățile vizuale sunt implementate ca
JavaScript nativ, inclus în blocul \texttt{<script>} din \texttt{app.html}.
Acesta rulează la fiecare încărcare de pagină și inițializează:

\begin{itemize}
  \item \textbf{Navbar scroll} --- eveniment \texttt{scroll} cu clasa
        \texttt{scrolled} la derulare
  \item \textbf{Hero visible} --- \texttt{setTimeout(100ms)} pentru animația
        de intrare a secțiunii hero
  \item \textbf{Reveal observer} --- \texttt{IntersectionObserver} pentru
        elementele cu clasa \texttt{.reveal}
  \item \textbf{ImageTilt} --- efect 3D pentru elementele \texttt{.tilt-container}
  \item \textbf{Tilt pe liste} --- efect 3D similar pentru \texttt{.project-images}
  \item \textbf{Toggle imagini} --- butonul \texttt{.img-thumb} comută între
        imaginea de exterior și cea de interior pe pagina Proiecte
\end{itemize}

% ═══════════════════════════════════════════════════════════════════════════════
%  CHAPTER 6 — BIBLIOGRAFIE
% ═══════════════════════════════════════════════════════════════════════════════
\chapter{Bibliografie}

\begin{itemize}[label=\textbullet]
  \item SvelteKit --- Documentație oficială: \url{https://svelte.dev/docs/kit}
  \item Svelte 5 Runes --- Documentație oficială: \url{https://svelte.dev/docs/svelte/runes}
  \item TypeScript --- Documentație oficială: \url{https://www.typescriptlang.org/docs/}
  \item Vite --- Documentație oficială: \url{https://vite.dev/guide/}
  \item MDN Web Docs --- CSS custom properties: \url{https://developer.mozilla.org/en-US/docs/Web/CSS/Using_CSS_custom_properties}
  \item MDN Web Docs --- IntersectionObserver API: \url{https://developer.mozilla.org/en-US/docs/Web/API/Intersection_Observer_API}
  \item MDN Web Docs --- requestAnimationFrame: \url{https://developer.mozilla.org/en-US/docs/Web/API/Window/requestAnimationFrame}
  \item @sveltejs/adapter-static --- Documentație: \url{https://kit.svelte.dev/docs/adapter-static}
  \item Frank Lloyd Wright Foundation --- Site oficial: \url{https://franklloydwright.org/}
  \item Wikipedia --- Frank Lloyd Wright: \url{https://ro.wikipedia.org/wiki/Frank_Lloyd_Wright}
  \item Space Grotesk (font) --- Google Fonts: \url{https://fonts.google.com/specimen/Space+Grotesk}
  \item DM Sans (font) --- Google Fonts: \url{https://fonts.google.com/specimen/DM+Sans}
\end{itemize}

\end{document}
